\begin{frame}{탐욕 정책: Q표에서 바로 행동이 읽힌다}
  $Q$를 $V$ 없이, $Q$만으로 쓰면:
  \[
  Q^{\pi}(s,a) = R(s,a) + \gamma \sum_{s'} P(s'\mid s,a)
  \sum_{a'} \pi(a'\mid s')\, Q^{\pi}(s',a')
  \]
  \vskip0.6em
  \begin{itemize}
    \item ``어떤 행동을 고를지''도 Q표에서 바로 --- \kb{탐욕 정책}
      (greedy policy):
      \[
      \pi(s) = \arg\max_a Q(s,a)
      \]
    \item $Q$를 학습한다 = ``가치 평가''와 ``행동 선택''을 \textbf{한 표로}
      모두 해결.
  \end{itemize}
  \vskip0.6em
  \begin{block}{작은 예 --- Q표의 모양을 잡자}
    상태 셋 장난감 MDP: $A$(출발)·$B$·$G$(도착, 터미널).
    $A$에서 \texttt{right} $\to$ $-1$로 $B$, \texttt{down} $\to$ 절벽 $-100$
    받고 $A$로. $B$에서 \texttt{right} $\to$ $G$(보상 0, 종료). $\gamma=1$.
  \end{block}
\end{frame}
